\scanpage{202}

% B01 / text / top, previous description
La description précédente de la notion de graphe a
l'avantage principal d'être ``autoduale'', \uline{et} de
séparer deux composantes de la structure :

% B02 / display / upper, (22), parts a and b
\[
 \begin{aligned}
 &\text{a) La }\unsure{\text{données}|\text{donnée}}\quad(S,S^*,R_*\subset S\times S^*)\\
 &\text{b) La donnée}\qquad \varphi_0:S\xrightarrow{\sim}S^*\ .
 \end{aligned}\tag*{(22)}
\]

% B03 / text / upper middle, without separation
Rien n'empêche d'ailleurs d'étudier de telles
structures, sans supposer que la relation a)
soit ``biséparante'' -- i.e. en ne la soumettant
à aucune hypothèse particulière. La \inserted{auto}dualité
(21) reste valable.

% B04 / text / middle, new question
Le jeu nouveau, c'est qu'on peut partir d'une structure a)

% B05 / display / middle, (23)
\[
 (S,S^*,R_*\subset S\times S^*), \tag*{(23)}
\]

% B06 / text / middle to lower, pregraph
sans donnée d'un $\varphi_0$, -- donc une structure
en un sens plus faible qu'une structure de
graphe, on pourrait l'appeler structure de
``pré-graphe'' -- structure considérée comme ``préliminaire''
à celle de graphe (séparé strict). Il se pose alors
la question d'étudier toutes les structures
de graphe correspondantes i.e. \uline{toutes} les
applications \unsure{bijections|bijectives} \deleted{(ou \unsure{$R_0$|$\mathcal R_0$}}

% B07 / display / lower, (24), left brace
\[
 \left\{\begin{aligned}
 &\varphi:S\xrightarrow{\sim}S^*\\
 &\text{telles que }\varphi\text{ soit ``}R_*\text{-\unsure{symétriques|symétrique} et antiréflexive''}\\
 &\quad\text{(on dit que }\varphi\text{ est une }R_*\text{-jauge)}
 \end{aligned}\right.\quad ; \tag*{(24)}
\]
% slanted
\margin{antiréflexivité n'est pas essentielle, si on admet des graphes avec boucles}

% B08 / text / foot, gauges and origin
\deleted{i.e. les}\inserted{l'ens des} ``$R_*$-jauges''. On s'intéresse surtout au cas,
bien sûr, où cet ensemble n'est pas vide -- et le $\varphi_0$
de tantôt jouera le rôle d'une ``origine'' dans l'ens

% archive sheet number: 201
